
\section{Elephant eye combinatorics}

We will use \cite{Book} for reference.

\sss{Description of the combinatorics}

Let $f: U' \ra U$ be an even  quadratic-like map with connected Julia set
$\Jul\equiv \Jul_f$ 
and both fixed points repelling.
Then one of these fixed point, called $\beta$, is the landing point of
the external $0$-ray,
while the other one, called $\alpha$, is the landing point of $\qq>1$
external rays $\Ray_i$ that a cyclically permuted with
some combinatorial rotation number $\pp/\qq$.
These rays partition $U$ into $\qq$ sectors $\Sec_i$ attached to $\alpha$,
counted anti-clockwise beginning with $S_0\ni 0$. 
See  \cite[\S 24.4]{Book}.


The symmetric rays $-\Ray_i$ land at the preimage $-\alpha$ of
$\alpha$.  Together with the rays $\Ray_i$, they partition $U$ into $2\qq-1$ domains: the
central one $\Ups$, the above $\qq-1$ sectors $\Sec_i$ attached to $\alpha$,
and symmetric sectors $\Zec_i$ attached to $-\alpha$.   
The intersections of these sets with $U'$ will be marked with
``prime'': $\Ups_i':= \Ups\cap U'$, etc.   

Assume that the map $f$ is renormalizable with period $p$ and little
(filled) Julia sets $K_n$, $n=0,1,\dots, p-1$, where $K_0\ni 0$
(see \cite[\S 28.4]{Book}).  We let $\Kfilled:= \bigcup K_n$.  

\msk
For the class of {\em elephant eye combinatorics} we assume:

\ssk\nin E1. The combinatorial rotation number of $\alpha$ is $1/\qq$.
  Then the domain $\Ups'$ is mapped onto $\Sec_1$ as a double covering,
  each $\Sec_i'$ and $\Zec_i'$ is univalently  mapped onto
  $\Sec_{i+1}$ for $i=1,\dots, \qq-2$, while each $\Sec_{\qq-1}$ and $\Zec_{\qq-1}$ is
  univalently mapped onto $\Ups \cup \bigcup \Zec_i\cup (-\Ray_i)$.

  
  \ssk\nin E2.
  $K_0$ is the only little Julia set contained in $\Ups$.

  \ssk\nin E3. $p= \qq+ \bb$, where $\bb$ is bounded by some $\bar \bb$
  (but $\qq$ is unbounded!).   We will also assume for simplicity that
    $\bb< \qq$. 

    \ssk Note that under these assumptions,
     each sector $S_i$, $i=1,\dots, \qq-\bb$,
    contains a single little Julia set, $K_i$.
    
  \sss{Hubbard tree}

  An efficient  way of describing the renormalization combinatorics is
  in terms of the Hubbard tree $\Hub\equiv \Hub_f$ (see \cite[]{Book}).
  Roughly speaking, it is the tree spanned by  the little Julia sets
  $K_i$ (collapsed to points) 
  inside the  big one.
  The vertices of the $\Hub$ are the collapsed $K_i$
(to which we will still refer as ``$K_i$''), together with the
$\alpha$-fixed point   and  all branched points.
(Note that for $\qq>2$, $\alpha$ is also a branched point.)%  
\footnote{The subtlety of the precise definition  has to do with the
  possibility  that $\Kfilled  $ can be non-locally-connected.}


\comm{*******
  Let $\Hub_i$, $i\in \Z/\qq\Z $ be the limbs of $\Hub $ attached to
  the $\alpha$-fixed point. We also let

  \ssk\nin $\bullet$
  $[K_i,K_j]\subset \Hub$
 be the arc of the tree connecting $K_i$ to $K_j$
(more generally, we can consider the span $[T_1, T_2, \dots, T_k]$ 
of any family of subtrees);

  \ssk\nin $\bullet$
  $\Hub^\bullet\subset \Hub_0$ be  
  the span of  $-\alpha$ and  all little Julia sets contained in
  $\bigcup \Zec_i$;
%   limb attached to $0$ and containing $\-\alpha$;

  \ssk \nin $\bullet$
    $\Hub^\#\subset \Hub$ be the span of $\alpha$ and all $K_i$ with $i>\qq$ contained in
    $\bigcup \Sec_i$. 

  Under the Elephant eye combinatorics, we have: 

 \ssk \nin $\bullet$    
    for $i\not=0$,  $\Hub_i$ is   homeomorphically mapped onto $\Hub_{i+1}$;

\ssk\nin 
$\Hub^\bullet$ is  homeomorphically mapped onto
****************}

\bignote{An interesting combinatorial lemma is hidden here} 

\comm{*********
\begin{lem}
  For $0< n< p$, let $\gamma_n$ be the arc of $\Hub$ connecting
  $K_n$ to either $\alpha$ or $-\alpha$ (depending on whether
  $K_n\subset \bigcup \Sec_i$ or $K_n\subset \bigcup \Zec_i$). 
If $K_m\in \gamma_n $, $0< m<p$, then 
 under the Elephant eye assumption, $m>n$.  
\end{lem}

\begin{proof}
  Let us consider the maximal   $n$   for which the statement is
  violated.

  Assume first that $n< p-1$ and
  $K_n \not\subset  \Sec_{\qq-1}\cup \Zec_{\qq-1} $.
Then the arc $\gamma_n$ is homeomorphically mapped onto
$\gamma_{n+1}\ni K_{m+1}$,
with  $m+1< n+1< p$,
contradicting the maximality property of $n$. 

\bignote{OK not to  distinguish notation for a set $K_n$ from
  a vertex $K_n$?} 

If $n=p-1$, then the arc $\gamma_n$ is homeomorphically mapped onto
  the arc $[\alpha, K_0]$ containing $K_{m+1}$.
  On the other hand,  by Assumption E2,  $[\alpha, K_0]$ does not
  contain any little Julia sets inside. 

Finally, assume $n< p-1$ but $K_n \subset \Sec_{\qq-1} \cup \Zec_{\qq-1}$.
All the more,  $m< p-1$, so by Assumption E2,  $K_{m+1}$ must be contained
in some lateral sector $\Zec_j$.
Then the arc $\gamma_n$ is homeomorphically mapped onto an arc
$\de\subset \Hub$ connecting $\alpha$ to $-\alpha$,
entering $\Zec_j$, and then proceeding to
$K_{m+1}$ and further  to $K_n$, without ever exiting $\Zec_j$. 
It follows that its piece from $-\alpha$ to  $K_{n+1}$ is equal to $\gamma_{n+1}$,
and moreover, it contains  $ K_{m+1}$.
This is a  contradiction once again. 
\end{proof}
***************************}

For $ n\not= 0\ \mod \, p$, let $\gamma_n$ be the arc of $\Hub$ connecting
  $K_n$ to either $\alpha$ or $-\alpha$ (depending on whether
  $K_n\subset \bigcup \Sec_i$ or $K_n\subset \bigcup \Zec_i$).%
\footnote{We tacitly assume E2, though the discussion can be easily
 extended  to the general case.} 
We also let $\gamma_0^+$ be the arc of $\Hub$ connecting $K_0$ to
$\alpha$, while  $\gamma_0^-$ be the arc  connecting $K_0$ to
$-\alpha$. We let $\gamma_0$ be the concatenation of these two,
i.e., the arc connecting $\alpha$ to $\alpha_-$.
We have the following transformation scheme for these arcs:

\ssk\nin $\bullet$
   Each $\gamma_0^\pm$ is homeomorphically  mapped by $f$ onto $\gamma_1$;
  
   \ssk\nin $\bullet$
   If $0< n< p-1$ and $K_n \not\subset  \Sec_{\qq-1}\cup \Zec_{\qq-1}$,
   then  $\gamma_n $ is homeomorphically  mapped by $f$ onto $\gamma_{n+1}$;

\ssk\nin $\bullet$
    If $0< n< p-1$ and $K_n \subset  \Sec_{\qq-1}\cup \Zec_{\qq-1}$,
   then  $\gamma_n $ is homeomorphically  mapped by $f$ onto
   $\gamma_{n+1}\cup \gamma_0$; 

 If $n = p-1$ (and hence $K_n \subset  \Sec_{\qq-1}\cup \Zec_{\qq-1}$
   then  $\gamma_n $ is homeomorphically  mapped by $f$ onto
   $ \gamma_0^+$.


\begin{lem}\label{bounded deg}
  Under the Elephant eye assumption,
  any edge $\gamma$ of  $\Hub$ is mapped under $f^p$ at most two-to-one
 into the union of at most $\bb+ 1$ consecutive  edges.
\end{lem}

\begin{proof}
It is sufficient to show this for the arcs $\gamma_n$, $0< n < p $, and for
$\gamma_0^\pm$, which  follows from the following observations:

\ssk\nin $\bullet$
Each arc $\gamma_0^\pm$ and each arc $\gamma_n$,
$1\leq n  \leq  \qq-2 $,  is mapped  under $f^{\qq-n-1}$
homeomorphically onto $\gamma_{\qq-1}$;

\ssk\nin $\bullet$
By induction, for $\qq-1 \leq n \leq  p-1 $,  and $k\le \bb$,
the arc $\gamma_n$   is mapped at most two-to-one by
$f^k$  into $\gamma_{n+k}\cup  \displaystyle { \bigcup_{0 \leq i\leq
    k-1}\gamma_i }$, and therefore

\ssk\nin $\bullet$
For $\qq-1 \leq n \leq  p-1 $,   the arc $\gamma_n$   is mapped under
$f^{\bb} $ at most two-to-one  into $\displaystyle { \bigcup_{0 \leq i\leq   \bb}\gamma_i }$.
\end{proof}

\subsubsection{Translation domains}

\bignote{Define them here?}

Let us consider a topological disk $\Om $  inside the translation region
which is the union of $2N+1$ consecutive sectors
$\Sec_n$, $1\leq  l-N \leq  n\leq l+N\leq \qq-1$.
Ler $\Om'$ be its pullback of $\Om_l$ under $f^p$ based at $K_l$,
i.e., the connected component of $f^{-p} (\Om)$ containing $K$. 
Let us also consider the arcs $\de_n\subset \Om$
connecting $\KK_n$ to $\KK_{n+1}$ and aligned with the Hubbard tree
(that is,  homotopic to the  concatenation of the edges $\gamma_n$ and
$\gamma_{n+1}$ considered  above). 

\begin{lem}\label{J-sets don't pull off}
 % There exsists $\bb'$ (depending on $\bb$ only) such that
Under the above circumstances, $\Om'$ contains the little Julia sets $K_n$,
 $   l- N + \bb \leq n\leq l+ N-\bb'$, and the arcs $\de_n$ connecting
 them.  
\end{lem}

\begin{proof}
Assume inductively in $k=0,1,\dots$ that the assertion is true for all
$\KK_n$ on combinatorial distance at most  $k$ from $\KK_l$.
Let for definiteness $n=l+k \geq l$ (and $n<\qq- \bb$). 
By Lemma \ref{bounded deg},  the  image $\om_n= f^p(\de_n) $ is contained in the
union of arcs $\de_{n+j}$, $  |j|\leq \bb$, so it is contsined in
  $\Om$.  It follows that $\de_n$ is the lift of $\om_n\subset \Om$
  based at $\KK_n$. Hence $\de_n$ is contained in $\Om'$, and so is $\KK_{n+1}$.  
\end{proof}



\begin{lem}\label{interesting moments}
 Let $(f^n)^* \Om\not\supset K_1$  %  $n\leq 3p$,
 be a pullback  of $\Om$
 containing a horizontal arc $\alpha$
 with endpoints on combinatorial distance at
most $N-\bb$ from  the base set $\KK_{l-N}$.
Then   $f^* \alpha$  does not pull off.
\end{lem}
 
\begin{proof}
Let the arc $\alpha$ lands on $K_i$ and $K_j$. By Lemma \ref{J-sets
  don't pull off}, none of these  Julia sets get lost under the pullback
(i.e.,  both of them are contained in $(f^{n+1})^*\, \Om$).
Let $\beta= f^* \alpha$ be the lift of $\alpha$ that begins on $K_{i-1}$. 
If it pulls off then  it ends on $K_{j-1}'$,
implying$K_{j-1}'\subset ( f^{n+1})^*  \Om$. 

Since  $( f^n)^* \Om $ does not contain $K_1$, 
the map $f: ( f^{n+1})^*\, \Om \ra ( f^n)^*\,  \Om $ is univalent. 
It follows that $K_{j-1}\not\subset ( f^{n+1})^*\,  \Om $ -- contradiction. 
\end{proof}
